% !PRINTED 46
\weberpage{46}
%<<CONT
\weq{11}{\begin{vmatrix} 1, & x_1, & x_1^2, & \sqrt{f(x_1)} \\ 1, & x_2, & x_2^2, & \sqrt{f(x_2)} \\ 1, & x_3, & x_3^2, & \sqrt{f(x_3)} \\ 1, & x_4, & x_4^2, & \sqrt{f(x_4)} \end{vmatrix} = 0.}

Aus dieser Gleichung kann man $x_1$ und $\sqrt{f(x_1)}$ \gesperrt{rational}
durch $x_2, \sqrt{f(x_2)}$; $x_3, \sqrt{f(x_3)}$; $x_4, \sqrt{f(x_4)}$ ausdrücken; denn
$x_1$ ist die Wurzel einer biquadratischen Gleichung, deren drei
andere Wurzeln $x_2, x_3, x_4$ sind, und deren Koeffizienten rational
von $x_2, \sqrt{f(x_2)}$; $x_3, \sqrt{f(x_3)}$; $x_4, \sqrt{f(x_4)}$ abhängen, und $\sqrt{f(x_1)}$ ist
mittels (11) rational durch $x_1$ darstellbar. Setzt man $x_4$ einer
beliebigen Konstanten gleich, so geht die Differentialgleichung
(10) in (1) über und (11) ergibt ihr Integral.

Wir wenden dies Theorem zunächst auf die Normalform an
und setzen
\[ f(x) = x(1-x)(1-\varkappa^2 x). \]
Nehmen wir in (11) $x_4 = 0$ an, so reduziert sich diese
Gleichung durch Division mit $\sqrt{x_1 x_2 x_3}$ auf:
\weq{12}{\begin{vmatrix} \sqrt{x_1}, & x_1\sqrt{x_1}, & \sqrt{(1-x_1)(1-\varkappa^2 x_1)} \\ \sqrt{x_2}, & x_2\sqrt{x_2}, & \sqrt{(1-x_2)(1-\varkappa^2 x_2)} \\ \sqrt{x_3}, & x_3\sqrt{x_3}, & \sqrt{(1-x_3)(1-\varkappa^2 x_3)} \end{vmatrix} = 0,}
oder wenn wir die Determinante nach den Elementen der ersten
Zeile entwickeln:
\weq{13}{\sqrt{x_1}\,(a + bx_1) + c\sqrt{(1-x_1)(1-\varkappa^2 x_1)} = 0,}
worin zur Abkürzung gesetzt ist:
\weq{14}{\begin{aligned}
a &= x_2\sqrt{x_2}\sqrt{(1-x_3)(1-\varkappa^2 x_3)} - x_3\sqrt{x_3}\sqrt{(1-x_2)(1-\varkappa^2 x_2)}, \\
b &= -\sqrt{x_2}\sqrt{(1-x_3)(1-\varkappa^2 x_3)} + \sqrt{x_3}\sqrt{(1-x_2)(1-\varkappa^2 x_2)}, \\
c &= -(x_2-x_3)\sqrt{x_2}\sqrt{x_3}.
\end{aligned}}

Wenn wir (13) rational machen, so erhalten wir die kubische
Gleichung:
\[ x(a+bx)^2 - c^2(1-x)(1-\varkappa^2 x) = 0, \]
deren drei Wurzeln $x_1, x_2, x_3$ sind. Es ist demnach identisch
\weq{15}{b^2(x-x_1)(x-x_2)(x-x_3) = x(a+bx)^2 - c^2(1-x)(1-\varkappa^2 x),}
%CONT>>
